\chapter{Introducción}

\section{Motivación y Contexto}

En los últimos años, el reconocimiento automático de acciones humanas en video ha emergido como un área clave dentro del campo de la visión por computadora, con aplicaciones que abarcan desde la seguridad y el entretenimiento hasta la medicina y el deporte. Más allá de identificar qué acción se realiza en un video, ha cobrado relevancia una tarea más compleja: la \textbf{Evaluación de la Calidad de Acciones} (\textit{Action Quality Assessment}, \ac{AQA}). Esta disciplina busca medir la calidad con la que se ejecuta una acción, considerando criterios como la técnica, la fluidez y la precisión del movimiento. A diferencia del reconocimiento de acciones, \ac{AQA} introduce desafíos adicionales tanto en términos de modelado como de evaluación.

El interés por \ac{AQA} ha crecido sustancialmente debido a su potencial impacto en contextos donde el desempeño técnico es crítico. En el ámbito deportivo, entrenadores y atletas pueden aprovechar sistemas automáticos para recibir retroalimentación objetiva y detallada que complemente la evaluación humana. En rehabilitación física, los profesionales podrían beneficiarse de modelos que evalúen en tiempo real la correcta ejecución de ejercicios terapéuticos, facilitando el seguimiento remoto de pacientes. En educación médica, la evaluación automática de gestos quirúrgicos podría apoyar la formación de residentes.

A pesar de estos beneficios potenciales, uno de los principales obstáculos para la adopción de sistemas de \ac{AQA} en escenarios reales radica en su elevada demanda computacional. Modelos de alto rendimiento basados en arquitecturas tridimensionales profundas, como I3D \cite{Carreira2017I3D} y SlowFast \cite{feichtenhofer2019slowfast}, han demostrado alta precisión al capturar patrones espacio-temporales complejos. Sin embargo, su tamaño y costo computacional los hacen inviables para aplicaciones que requieren baja latencia, eficiencia energética o procesamiento local en dispositivos con recursos limitados.

\section{Planteamiento del Problema}

La estimación automática de la calidad de acciones en video requiere capturar información espacial y temporal asociada a la ejecución de una acción humana. Los modelos 3D profundos, como \ac{I3D}, han sido ampliamente utilizados por su capacidad para modelar dependencias espacio-temporales complejas; sin embargo, su alto costo computacional limita su uso en escenarios donde se requiere baja latencia, menor consumo de memoria o despliegue en dispositivos con recursos limitados.

La literatura reciente ha propuesto la destilación de conocimiento (\textit{Knowledge Distillation}, \ac{KD}) \cite{Hinton2015Distilling} como estrategia para transferir información desde modelos \textit{Teacher} pesados hacia modelos \textit{Student} livianos. Esta línea de trabajo \cite{Yu2021GroupAware,Mun2019TKD,Cai2023VisionLanguageAQA} parte de la premisa de que las arquitecturas livianas presentan una brecha sustancial de rendimiento frente al \textit{Teacher} y que dicha brecha debe cerrarse mediante señales adicionales de supervisión. No obstante, esta premisa fue formulada en un contexto donde los pesos preentrenados disponibles, las prácticas de entrenamiento y las arquitecturas eficientes eran menos maduras.

Actualmente, los modelos livianos pueden beneficiarse de pesos modernos preentrenados en ImageNet, optimizadores como AdamW, programación cosenoidal de la tasa de aprendizaje, entrenamiento con precisión mixta, acumulación de gradientes y módulos temporales eficientes como \textit{Temporal Shift Module} \cite{Lin2019TSM}. Bajo este nuevo escenario, no está claro si los \textit{Students} livianos siguen requiriendo destilación para aproximarse al rendimiento de un \textit{Teacher} 3D, ni cuáles son los límites reales de un \textit{pipeline} liviano moderno para \ac{AQA}.

A partir de esta problemática, surge la siguiente \textbf{pregunta de investigación}:

\begin{quote}
\textbf{¿En qué medida un \textit{pipeline} liviano basado en arquitecturas modernas (TSM-MobileNetV2 y MobileNetV3), correctamente inicializado y entrenado con prácticas actuales, puede aproximarse al rendimiento de un \textit{Teacher} 3D profundo en \ac{AQA} a un costo computacional sustancialmente menor, y dónde están sus límites de aplicabilidad?}
\end{quote}

Este trabajo aborda dicho problema mediante el desarrollo de un \textit{pipeline} eficiente basado en arquitecturas livianas modernas, evaluando su capacidad para aproximarse al rendimiento de \ac{I3D} y \textit{SlowFast} en distintos dominios, caracterizando su costo computacional y delimitando los escenarios donde el enfoque resulta suficiente o requiere mecanismos adicionales.

\section{Objetivos}

\subsection{Objetivo General}

Desarrollar un \textit{pipeline} eficiente de \textit{Action Quality Assessment} basado en arquitecturas livianas modernas, evaluando empíricamente su capacidad para alcanzar rendimiento competitivo respecto a modelos 3D profundos en dominios de naturaleza distinta y delimitando sus límites de aplicabilidad.

\subsection{Objetivos Específicos}

\begin{itemize}
  \item \textbf{OE1.} Revisar el estado del arte en \ac{AQA} y métodos de compresión o destilación aplicados a video, identificando las principales suposiciones sobre la brecha entre modelos \textit{Teacher} 3D y \textit{Students} livianos.

  \item \textbf{OE2.} Construir un \textit{pipeline} reproducible de entrenamiento e inferencia para \textit{Students} livianos, específicamente \ac{TSM}-MobileNetV2 y MobileNetV3, empleando pesos preentrenados, optimización moderna, precisión mixta y acumulación de gradientes.

  \item \textbf{OE3.} Cuantificar la brecha de rendimiento entre los \textit{Students} livianos y el \textit{Teacher} \ac{I3D} mediante \ac{SRCC}, \ac{PLCC} y \ac{MAE} en datasets representativos de \ac{AQA}, complementando la comparación con un \textit{Teacher} 3D más moderno (SlowFast-R50) en el dataset principal.

  \item \textbf{OE4.} Analizar el impacto de componentes clave del \textit{pipeline} mediante ablaciones mínimas sobre el dataset principal, priorizando las decisiones con mayor efecto esperado en el rendimiento (preentrenamiento ImageNet y módulo temporal).

  \item \textbf{OE5.} Caracterizar la eficiencia computacional del \textit{pipeline} propuesto en términos de parámetros, \ac{FLOPs} y latencia de inferencia frente al \textit{Teacher} \ac{I3D}.

  \item \textbf{OE6.} Delimitar los límites de aplicabilidad del \textit{pipeline} mediante análisis complementarios de transferencia cruzada entre dominios y aporte marginal de la destilación de conocimiento sobre la línea base liviana.
\end{itemize}

\section{Organización de la Tesis}
\label{sec:organizacion}

El presente documento está estructurado de la siguiente manera:

\begin{itemize}
    \item \textbf{Capítulo 2: Marco Teórico} -- Fundamentos de \textit{Action Quality Assessment} (\ac{AQA}), métricas clave (\ac{SRCC}, \ac{PLCC}, \ac{MAE}), conjuntos de datos (AQA-7, MTL-AQA, JIGSAWS) y nociones de \ac{KD}.

    \item \textbf{Capítulo 3: Estado del Arte} -- Revisión de técnicas recientes en \ac{AQA} y \ac{KD} aplicada a video, con énfasis en las suposiciones sobre la brecha entre \textit{Teacher} y \textit{Student}.

    \item \textbf{Capítulo 4: Propuesta} -- Metodología completa del \textit{pipeline} liviano: arquitecturas \textit{Student} (\ac{TSM}-MobileNetV2 y MobileNetV3), receta de entrenamiento moderna, datasets y protocolo de evaluación. Se incluye también el módulo de destilación que se utiliza como análisis marginal en el Capítulo 5.

    \item \textbf{Capítulo 5: Resultados} -- Comparación principal del \textit{pipeline} liviano frente a los \textit{Teachers} \ac{I3D} y SlowFast-R50, ablaciones sobre el preentrenamiento ImageNet y el módulo temporal, eficiencia computacional, transferencia cruzada y análisis marginal de la destilación.

    \item \textbf{Capítulo 6: Conclusiones} -- Hallazgos clave alineados a los objetivos específicos, limitaciones y direcciones futuras.
\end{itemize}
